\section{Introduction}

An investor restricted to listed options has a simple question with a difficult answer:
which calls and puts actually deserve premium capital? The usual asset-allocation template
starts from return columns. Options do not fit that template cleanly. They are funded,
expiring, state-contingent cashflows whose economics depend on premium paid or received,
strike, expiry, exercise style, settlement, volatility, convexity, and the state of the
underlying.

This paper builds a premium-weighted option-only Markowitz framework for that choice. The
universe can contain arbitrary listed calls and puts across equities, strikes, expiries, and
moneyness buckets. Positions can be long or short. Cash is not optimized as a separate
risky asset; it is the NAV numeraire and collateral account used to scale option premium
exposures. The premium is still paid or received. A long option that expires worthless loses
its premium weight, and a short option that finishes in the money pays the payoff through
the same NAV accounting. The allocation problem is therefore not ``which stock return
column has the best mean,'' but whether a funded option cashflow deserves capital after
settlement, costs, beta exposure, Greek budgets, and validation gates.

\subsection{Research Question and Contributions}

The research question is: how should an allocator choose among listed calls and puts when
every instrument is a funded, expiring, state-contingent cashflow rather than a standard
asset return column? The proposed answer is a premium-weighted option-only Markowitz
framework that maps listed options into NAV-normalized cashflows, payoff and settlement
conventions, Greeks, residual risk, conditional option-risk-premium forecasts,
implementation screens, and portfolio constraints. In this sample and under these
conventions, the validated empirical claim is that the framework is coherent,
point-in-time, auditable, and useful for diagnosing whether option convexity, VIX exposure,
or short-premium risk deserves capital. The non-claim is equally important: the article does
not claim live alpha, production tradability, broker-executed performance, live margin
parity, or deployable option trading performance.

The organizing thesis is that option-only portfolio construction should be written as
conditional option cashflow engineering. A listed option is not a stock with a different
ticker. It is a state-contingent claim with delta, gamma, vega, theta, exercise, settlement,
liquidity, and marking conventions. Two options can share the same underlying and still
contribute very different risks. A put and a call can leave the book with little net delta and
large volatility or convexity exposure. The covariance matrix must therefore come from the
instrument map, not only from historical option-price correlations.

The theoretical contribution is the option-only analogue of Markowitz. Classical portfolio
theory starts with expected returns and an asset covariance matrix
\citep{markowitz1952portfolio}. This paper starts with option premium weights, a local
Greek return map, conditional option-risk-premium forecasts, and a residual covariance term.
Once those objects are built, the mean--variance algebra is familiar: the unconstrained
maximum-Sharpe direction is proportional to
\(\widehat\Sigma_{O,t}^{-1}\widehat\mu_t\). The difference is the construction of
\(\widehat\Sigma_{O,t}\) and \(\widehat\mu_t\), not the final quadratic geometry.

What is new is the object rather than a single trade. The paper is not a defense of long
calls, short puts, covered calls, protective puts, or a named volatility spread. It asks how
to allocate across an arbitrary option universe once each contract has a mark, a payoff
convention, a Greek exposure vector, a conditional expected return, and an implementation
ledger. The expected-return side is also explicit: covariance optimization alone cannot
create a high Sharpe ratio. A useful option allocator needs forecastable premia such as
carry, variance risk premia, volatility risk premia, skew compensation, tail-insurance
pricing, relative value, and regime-dependent convexity.

\subsection{Related Work and Positioning}

The starting point is mean--variance allocation. \citet{markowitz1952portfolio} showed how
expected returns and a covariance matrix define an efficient frontier, while
\citet{sharpe1966} made reward-to-variability central to performance evaluation. Listed
options require a different first step because the columns are not ordinary asset returns.
Black--Scholes--Merton option valuation \citep{blackscholes1973,merton1973} and the
Black--76 forward option model \citep{black1976commodity} provide local pricing and Greek
maps, not a complete physical-measure portfolio allocation rule. The present framework uses
those maps to translate options into delta, gamma, vega, theta, VIX-forward exposure, and
residual risk before applying Markowitz geometry.

Expected option returns have their own economic sources. Empirical option-return research
links compensation to carry, volatility risk premia, skew, crash insurance, and tail-risk
states \citep{coval_shumway2001,bakshi_kapadia2003,carr_wu2009,
bollerslev_tauchen_zhou2009,broadie_chernov_johannes2009}. This article uses that spine to
motivate conditional option-risk-premium forecasts rather than treating the optimizer as a
source of alpha. Robust covariance estimation also matters because inverse-covariance
portfolio rules amplify noisy inputs; shrinkage follows \citet{ledoitwolf2004}. Evaluation
uses both reward-to-volatility and downside-aware measures, including Sortino
\citep{sortino_price1994}, Omega \citep{keating_shadwick2002}, Calmar and drawdown
\citep{magdon_ismail_atiya2004}, and active risk \citep{grinold_kahn2000}.

Implementation frictions and validation discipline determine the strength of any empirical
claim. Option overlays are sensitive to spread, settlement, borrow, assignment, margin, and
capacity assumptions. Backtest overfitting and multiple testing can turn optimizer variants
into false discoveries, so the analysis uses bootstrap inference and reality-check ideas
consistent with \citet{bailey2017pbo} and \citet{lopezdeprado2018}. The distinction from
strategy-specific option papers is therefore deliberate. The object is not a covered-call,
put-write, crash-put, or volatility-spread test. The object is an instrument-level cashflow
accounting and allocation framework for arbitrary listed options before asking whether any
option sleeve deserves capital.

The empirical result is intentionally mixed. The gross point-in-time backtest tests whether
the option-only Markowitz object behaves coherently on the local OPRA-derived panel. The
post-cost research layer subtracts conservative spreads, fees, slippage, borrow, capacity,
simulated margin funding, and assignment-risk penalties. VIX options are included as
volatility derivatives with VX-forward Black--76 Greeks, but VIX option evidence is
headline-grade only when exact VRO/SOQ settlement supports the relevant rows. If the
coverage table reports proxy rows, the same evidence is reported only as diagnostic
hedge-sleeve evidence rather than a listed-settlement claim.

The article proceeds in the same order as the research workflow. Section 2 defines the data
and empirical design. Section 3 builds the option cashflow and Greek-risk framework.
Section 4 describes the conditional portfolio construction model. Section 5 sets out the
metrics, point-in-time validation, inference, and claim gates. Section 6 reports the
empirical results and negative evidence. Section 7 concludes with the limits of the current
research implementation. Appendix A contains supplemental figures, option/VIX diagnostics,
robustness tables, instrument conventions, proof sketches, and reproducibility notes.
